\section{Introduction}
\label{sec:introduction}
\todo{Can be kept as is, I would say.  Review: Change wording in
  contributions to differentiate ourselves more from CIP /
  Kaleidoscope - it currently isn't very prominent that we get rid of
  constraint constructors}

Constraints and constraint programming occur in a variety of
application domains within computer science and related fields,
including graphics, CAD/CAM, planning and operations research,
artificial intelligence, user interface toolkits, and programming
languages. A \emph{constraint} in these domains is a relation that
should hold, for example: that there be a minimum of 10 pixels
horizontal space between two buttons on a screen, that a resistor in
an electrical circuit simulation obey Ohm's Law, that a maximum of 10
parts per hour can be produced by a machine in a factory, that all of
the digits in a row or column or subregion of a Sudoku puzzle must be
different, or that the height of a bar in a bar chart correspond to a
value produced by a simulation. Constraints are declarative: they
specify \emph{what} should be the case rather than \emph{how} to
achieve it.  They thus provide flexibility in the way they can be
employed, since a single constraint can typically be used in multiple
ways.

A number of features are desirable in constraint programming to solve
a useful set of problems.  Unfortunately it is in general much easier
to state constraint problems than to solve them --- some problems are
extremely difficult, and others are undecidable.  (For example,
suppose we have a constraint that $x^n+y^n=z^n$ for integers $x$, $y$,
and $z$, with $n>2$.  Waiting 358 years for the solver to determine
that the constraint is unsatisfiable is problematic.  Or as another
example, a language that allows the programmer to state 
the Post Correspondence Problem as a constraint
would thus allow undecidable constraints to
be specified.)  On the other hand, for some restricted but
very useful classes of constraints, quite efficient solvers are
available.  Therefore, practical solvers solvers impose restrictions
on the kinds of constraints they can solve, for example, by supporting
only floats and not integers or requiring numeric equalities to be
linear.

So, in general, a variety of solvers are
needed, depending on the type domain and particular characteristics of
the constraints.  It is also useful to have these solvers interoperate
to find a solution to a particular problem. Furthermore, in some
application domains, it is useful to support soft as well as hard
constraints, that is, constraints that should be satisfied if possible
and constraints that must be satisfied. Finally, in the presence of
frequent changes incremental solvers are required that can efficiently
re-solve a set of constraints.  Interactive graphical systems are an
example of all of these: they include constraints on complex objects
such as points and areas as well as floats and integers, soft
constraints to express preferences regarding the layout, and require
incremental solvers to efficiently re-solve constraints repeatedly as
a part of the figure is moved.

There are two main options how to use constraints in imperative
programs. The most common way is via a constraint satisfaction
library.  This means that no changes are required to the underlying
programming language, but has the disadvantage that, besides writing
the constraints, the solver needs to be invoked explicitly at
appropriate times during the execution. If it is not, the programmer
may inadvertently ignore or bypass the constraints.  An alternate
technique is to support constraints directly in the underlying
programming language.  This addresses the problem of inadvertently
bypassing the constraint system, because the programming language
ensures that the solver is invoked whenever neccessary. This approach
also typically provides a more convenient syntax for writing
constraints, either using a separate \ac{dsl} embedded in the host
language (with the drawback of having to learn essentially two
programming languages), or (as in the work reported here) by allowing
constraints to be written in the host language itself, of course at
the cost of having to support a new or extended programming
language.

Prior approaches of the latter kind typically provide a unified
runtime, a semantics for the interactions between declarative and
imperative code, and a linguistic symbiosis. As an extension of these,
this work additionally unifies the programming constructs for
encapsulation and abstraction by using only object-oriented methods,
classes, and inheritance for defining behavior and structure in both
the imperative and declarative paradigm.

We present our design for integrating constraints with a dynamic,
object-oriented language that preserves a familiar imperative
programming model, called \langname, and describe an implementation of
that design as an extension to Ruby~\cite{flanagan2008ruby}, called
\toolname. Key contributions for \langname with respect to language
design are as follows.

\begin{itemize}
\item The semantic model for \langname is a direct extension of the
  object-oriented model (with dynamic typing, object encapsulation,
  classes and instances or prototypes with methods, and message
  sends).  It supports placing constraints on the results of message
  sends rather than just on object attributes --- thus, we argue,
  being more compatible with object-oriented encapsulation and
  abstraction than prior approaches in \ac{cip}.
\item The syntax of \langname is a strict superset of that of the base
  language -- with only one minor extension to allow a question-mark
  as method name -- making it easy for the programmer to write and
  read constraints and object-oriented code.
\item \langname includes constructs that allow specifying soft as well
  as hard constraints, and that support incremental solving.
\end{itemize}

Additional contributions with respect to implementation techniques are
as follows.

\begin{itemize}
\item We present a technique for implementing object constraint
  languages, which uses a primitive to switch the interpreter between
  imperative evaluation, constraint construction, and constraint
  solving modes.  In imperative evaluation mode, the interpreter
  operates in the standard fashion, except that \load and \store
  operations on variables with constraints obtain the value from the
  constraint solver or send the new value to the solver (which may
  trigger a cascade of other changes to satisfy the constraints).  In
  constraint construction mode, the expression that defines the
  constraint is evaluated, not for its value, but rather to build up a
  network of primitive constraints that represent the constraint being
  added.  The interpreter keeps track of dependencies in the process,
  so that, as needed, the solver can be activated or the code to
  construct the constraint can be re-evaluated.
\item \langname provides an architecture for constraint solvers that
  supports multiple constraint solvers, which makes it straightforward
  to add new solvers, and that does not privilege the solvers provided
  with the basic implementation (they are simply the solvers that are
  in the initial library).
\item We describe a working prototype system, integrated with a state
  of the art Ruby virtual machine and \ac{jit} compiler. In the
  absence of constraints, the performance of a program written in the
  host language (Ruby) is only modestly impacted.
\end{itemize}

\toolname is a general-purpose object constraint language, and we have
implemented a variety of example programs, including a video streaming
example with both hard and soft constraints on performance,
interactive graphical layout examples (which exercise soft constraints
and incrementality), electrical circuit simulations, and solvers for
puzzles such as Eight Queens and cryptarithmetic.

\section{Background}
\label{sec:background-related-work}

Programs frequently have some set of constraints that should hold.
This could be that all parts in an electrical circuit simulation obey
the laws of physics, that the rows in a Sudoku include each digit from
0 to 9, or that a streamed video play smoothly in the presence of
changing CPU and network load.

\todo{We should start with the part about constraints in programs, and
  give a few examples. Show differences between one-way and multi-way
  constraints, constraint domains, give examples for solvers that can
  solve these.

  Review: Then make clear that we are interested not in enumerating
  solutions, but just in finding any solution - then reference
  self-healing/resilient systems/data-structure repair (Rinard
  et.al.)}

These constraints have different properties (domain, dynamicism,
one-way, multi-way). In this work, we only consider dynamic
constraints and system invariants to support writing resilient and
self-healing systems, not enumerating solutions to a constraint model.

For each of these, a number of different solvers may be used to solve
them...

For imperative languages, the usual approach to dealing with such
constraints is to leave it entirely up to the programmer to ensure
that they are satisfied --- the constraints are either just in
comments and documentation, or perhaps in the form of
machine-checkable assertions.  For the latter case, programmers
typically write assertions to fail early if these constraints are
unexpectedly not satisfied~\cite{rinard2005exploring}.

\section{Related Work}
The background section should have introduced the rectangle
constraints. Now give a solution in imperative terms for that. Give a
logic programming solution. Compare, mention again that we want to use
it in an imperative setting, an re-use the OO methods for area and
visibility.

Rewrite the solution using the pure-library approach.

Rewrite using a DSL

Rewrite using dataflow

Rewrite using CIP

Review: Rewrite using Ilog, Squander, or a similar
annotation/pre-post-condition/macro system Review: Look at Kaplan

Rewrite using OCP to introduce next section (?)

\section{Object Constraint Programming}
\label{sec:object-constraint-programming}
What is OCP? It is an integration of constraints with OO. It is
similar to CIP, but instead of adding separate constructs for
constraint-constructors and -structures, we re-use object-oriented
methods and types. The only two semantic extensions are for constraint
creation and solving - everything else is just convenience or library
code.

Review: we offer an informal description of our two semantic
extensions. This work is very recent, and we are still determining
details by using this in applications. A detailed, formal semantics
isn't useful at this point.

\subsection{Constraint Construction - the always primitive}
\label{sec:constr-constr-always}
This section should describe that OCP proposes the always method as a
mechanism similar to "eval", but that does not evaluate the code and
return the result, but instead constructs a constraint from the
expression in the code. In the process, all variables that are found
to contribute to the expression are marked as constraint variables,
and a representation for them is created that can be passed to a
solver

\subsubsection{Restrictions on Constraint Expressions}
\label{sec:restr-constr-expr}
No side-effects

\subsubsection{Constraints and Control Structures}
structures work, but only if all variables are in single-assignment
form (so loops with counters do not qualify)

\subsubsection{Constraints on Identity, Type, or Protocol}
Just explain that these can be checked

\subsubsection{Preferential Constraints}
Explain why we want them, and that they are solver specific.

Review: Ideally, solvers use the same language for preferences (high,
medium, low? Strong, medium, weak?), but it is ultimately up to the
library (and the arguments are passed to the library - they have no
meaning for the always-primitive)

Review: these are not required, but very useful

Review: Solvers are free to throw an exception when asked for a
preference constraint that they don't support (i.e. Z3 and not
required)
	
\subsection{Reading and Writing Constrained Variables}
This section should describe the semantics reading a constrained
variable (it is read from the solver instead), and writing a
constrained variable (either the new value is put into a solver, or
the expression is re-calculated)
	
\subsubsection{Stay Constraints}
\subsubsection{Read-only Variables}
A syntax for this is nice - we provide one for Babelsberg/R (ref Impl
section on Constrained Variables), but it is not necessary
\subsubsection{Incremental Resatisfaction}
This is needed because … (explain) It is, in this work, not
semantically significant, and can be implemented in library code. An
example of such an implementation is given in the implementation
section
			
Review: "edit" is not a language construct, just a library method

\subsection{Solving Constraints}
All constraints must be solved, but generally, solving is hard - we
propose that solvers are provided as libraries for easy
extensibility. The set of solvers should be chosen per application by
the programmer - just as he chooses other libraries appropriate for
the application. Solvers register themselves for specific types, and
are then automatically used for those. The programmer can also specify
a specific solver when calling always.
	
Solvers receive additional options passed to always (such as weight
and priority) - it is up to the solvers to throw an error. For
example, Z3 does not support preferential constraints, so it should
throw an error
	
Review: Solvers have to figure out how to map types. Z3 only supports
reals, but the C-API already maps these to C-doubles, so we just use
that interpretation (although there may be rounding errors, we just
don't care - if it doesn't work for you, use something else)

\subsection{Constraint Meta-Objects, controlling constraint duration
  and properties}
Explain what these are for - they provide access to enable/disable
constraints, change their strength, change their solver
	
\section{Implementation}
Review: Ruby is chosen because it was available and easily
extended. We have an implementation in Javascript, too.

\subsection{Constrained Variables}
Read-only constraint impl, stay constraint impl, syntactic extension

Review: GC

\subsection{Execution Contexts}
Review: this needs code examples

\subsection{Implementing Edit Constraints}

\subsection{Adding New Solvers}
Review: Array Solver
	
\section{Evaluation}
\subsection{Performance Evaluation}
Review: the JIT removes the overhead of accessing unconstrained
variables. The remaining overhead is an artifact of the code that is
not JITted

Review: we assume that most programs will not only execute constraints
- to rewrite a program in this way would mean to rewrite it in prolog
(or another logic programming language) - of course they can use the
new features extensively - but constraints ultimately don't "do" much
(in terms of user-visible side-effects like drawing to the screen or
writing to disk). Programs that do these things have other bottlenecks
than constraint solving performance.

Review: The difference between edit/non-edit is expected, because the
"planning" is basically re-run everytime

Review: the constraint version is more generic (we can tweak any
value, not just the mouse location), but lets assume this refactoring
did not try to make the code more general, just the constraints more
readable

\subsection{Comparison with other approaches}
Review: Array Solver benchmark

Review: Linked List (identity constraints) benchmark

Review: DeltaBlue example

Review: Video streaming app - less details, just report that solving
costs X time, and mention that overall time per frame is Y (with Y
much larger X)
	
\section{Conclusion and Future Work}


%%% Local Variables:
%%% mode: latex
%%% TeX-master: "babelsberg_jot"
%%% End:
